%%%%%%%%%%  scan idx 80  |  folio 67  %%%%%%%%%%
N.B.
Toutes les formules «~universelles~» liant $i_!$ et $i^!$ (resp.\ $i_*$ et $i^*$) sont conséquences formelles des précédentes, comme on peut s'en convaincre sans grande difficulté.

De la formule (v) résulte la formule
\[
 i_!(\gamma^p(N,x))=\gamma^p(i_!(x)).
\]
La conjonction du théorème 2.14 et de la proposition 1.5 montre que $\gamma^p(N,x)$ est, pour $p=q+j$, de filtration $\geq j$~; si l'on note $G_p(N,x)$ sa classe dans $G^j(\underline K(Y))$, le théorème 2.16 montre que l'on a
\[
 i_*(G_p(N,x))=c^p(i_!(x)).
\]
Enfin, la proposition 1.5 jointe au théorème 2.14 permet d'exprimer $G_p(N,x)$ comme un polynôme à coefficients entiers en les $c^i(N)$, les $c^i(x)$ et $\varepsilon(x)$~; on en déduit le théorème~:

\textbf{Théorème 2.17.} \emph{Supposons le corps $k$ de caractéristique $0$ et soit $i:Y\longrightarrow X$ un morphisme d'injection de variétés quasi-projectives non singulières. On a alors les formules (i) à (v), et de plus, pour tout entier $j\geq0$, la formule~:}
\begin{equation}\tag{2.37}
 c^{q+j}(i_!(x))=i_*\left(\frac{\widetilde c(x)*\lambda_{-1}(\widetilde c(N))}{(-1)^qN^q}\right)^{(j)}.
\end{equation}

Le deuxième membre de la formule (2.37) est une façon abrégée de désigner le polynôme en les composantes de $\widetilde c(x)$ et $\widetilde c(N)$ décrit dans la proposition 1.5.

\textbf{Corollaire :} \emph{Dans les conditions précédentes, on a~:}
\begin{equation}\tag{2.38}
 \ch(\widetilde c(i_!(x)))=i_*\left(\ch(\widetilde c(x))\,\mathcal T(\check N)^{-1}\right)
\end{equation}
\emph{où $\ch$ est l'homomorphisme de Chern défini par la formule} (1.26).

